\chapter{Lineære ligninger i lineær algebra}\label{chap:ch7}

I 1949 brugte Harvard-professor Wassily Leontief en af de tidligste computere, Mark II, til at løse et lineært ligningssystem, der modellerede den amerikanske økonomi. Han inddelte økonomien i 500 sektorer, fx kul, bilindustri og kommunikation, og beskrev med lineære ligninger, hvordan sektorerne påvirkede hinanden. På grund af hardwarebegrænsninger reducerede han problemet til 42 ligninger, som det tog Mark II 56 timer at løse. Arbejdet var en milepæl for anvendelsen af computere til beregninger og indbragte Leontief Nobelprisen i 1973.

Leontiefs arbejde viste, hvordan lineær algebra kombineret med computere kunne løse store problemer -- en idé, der har endnu større betydning i dag inden for softwareudvikling. Lineær algebra er grundlæggende for mange moderne teknologier, især maskinlæring (ML) og kunstig intelligens (AI). Fra anbefalingssystemer til optimering af neurale netværk er lineær algebra kernen i de algoritmer, der driver nutidens intelligente systemer.

I ML og AI omdannes store datasæt til matricer, hvor teknikker fra lineær algebra som matrixfaktorisering, egenværdidekomposition og singulærværdidekomposition (SVD) bruges til at udlede indsigter og foretage forudsigelser. Neurale netværk, som er rygraden i AI, bygger på lineær algebra, når data sendes gennem lagene, og når vægtene justeres under træningen. Som softwareingeniør får du ved at mestre disse begreber i lineær algebra de redskaber, der skal til for at bygge skalerbare, intelligente systemer, som kan håndtere nutidens mest komplekse beregningsproblemer.

Med den eksponentielle vækst i data og regnekraft bliver lineær algebra stadig vigtigere i softwareudvikling, især inden for AI og ML. Det er et afgørende værktøj, der bygger bro mellem teoretisk matematik og virkelige anvendelser inden for teknologi.

\section{Lineære ligningssystemer}

Lineære ligninger og ligningssystemer udgør grundlaget for lineær algebra, et centralt område af matematikken med omfattende anvendelser inden for softwareudvikling. Fra computergrafik til algoritmer til maskinlæring er det afgørende at forstå, hvordan man modellerer og løser lineære ligningssystemer, for at kunne udvikle effektive softwareløsninger. Vi begynder med en definition:

\begin{definition}{Lineære ligninger} En lineær ligning i variablene $x_1, x_2, \ldots, x_n$ er en ligning, der kan skrives på formen
\[
a_1 x_1 + a_2 x_2 + \cdots + a_n x_n = b
\]

hvor $b$ og koefficienterne $a_1, \ldots, a_n$ er konstanter.

\end{definition}

Lineære ligninger beskriver rette linjer, planer og hyperplaner i forskellige dimensioner, hvilket gør dem nyttige til at modellere sammenhænge i data og algoritmer.

\begin{example} En linje i det todimensionale rum givet ved $y = m x + b$ er en lineær ligning. Den kan omskrives, så den passer til vores standardform:

    \[
-m x+y=b
    \]

Her er $a_1 = -m$, $a_2 = 1$, og $b$ er konstantleddet. \end{example}

\begin{example} Den generelle ligning for en plan i det tredimensionale rum er

\[
a x+b y+c z=d
\]

hvor $a$, $b$, $c$ og $d$ er konstanter. Ligningen er lineær i variablene $x$, $y$ og $z$.
\end{example}

\begin{definition}{Lineære ligningssystemer} Et \textbf{lineært ligningssystem} (eller blot et \textbf{lineært system}) er en samling af en eller flere lineære ligninger i de samme variable. \end{definition}

\begin{example} Se på følgende lineære ligningssystem i variablene $x_1, x_2, x_3$:
\[
\begin{aligned}
2x_1 + 3x_2 + x_3 &= 3 \\
7x_2 - 4x_3 &= 10 \\
x_3 &= 1
\end{aligned}
\]
Systemet består af tre ligninger med tre ubekendte.

\end{example}

For at løse systemet bemærker vi først ud fra den tredje ligning, at $x_3 = 1$. Vi indsætter dette i den anden ligning og løser med hensyn til $x_2$:
\[
\begin{aligned}
7x_2 - 4(1) &= 10 \\
7x_2 &= 14 \\
x_2 &= 2
\end{aligned}
\]
Dernæst indsætter vi $x_2 = 2$ og $x_3 = 1$ i den første ligning og løser med hensyn til $x_1$:
\[
\begin{aligned}
2x_1 + 3(2) + 1 &= 3 \\
2x_1 &= -4 \\
x_1 &= -2
\end{aligned}
\]
Løsningen til systemet er altså $\left(x_1, x_2, x_3\right) = (-2, 2, 1)$.

\begin{definition}{Løsninger til et lineært ligningssystem} En \textbf{løsning} til et lineært ligningssystem i variablene $x_1, x_2, \ldots, x_n$ er en liste af tal $(s_1, s_2, \ldots, s_n)$, der opfylder alle systemets ligninger, når tallene indsættes på variablenes plads $x_1, x_2, \ldots, x_n$ i den rækkefølge. Mængden af alle mulige løsninger kaldes systemets \textbf{løsningsmængde}. To lineære ligningssystemer kaldes \textbf{ækvivalente}, hvis de har samme løsningsmængde. \end{definition}

Et lineært ligningssystem har et af følgende udfald:
\begin{enumerate}[label=(\roman*)]
    \item \textbf{Ingen løsninger}, når ligningerne er inkonsistente.
    \item \textbf{Præcis én løsning}, når der er ét entydigt sæt værdier, som opfylder alle ligningerne.
    \item \textbf{Uendeligt mange løsninger}, når der er flere sæt værdier, som opfylder ligningerne.
\end{enumerate}

Vi siger, at et system er \textbf{konsistent}, hvis det har enten én eller uendeligt mange løsninger. Vi siger, at et system
er \textbf{inkonsistent}, hvis det ikke har nogen løsning. Vi formaliserer dette senere i kapitlet.

\begin{remark} At afgøre, om et system er konsistent, handler om løsningers \emph{eksistens}. Hvis der findes løsninger, kan vi desuden undersøge, om løsningerne er \emph{entydige}. \end{remark}

\subsection*{Ligningssystemer skrevet med matricer}
Matricer giver en kompakt og effektiv måde at skrive og behandle lineære ligningssystemer på, hvilket især er en fordel i softwareanvendelser med store datasæt.

\begin{definition} En \textbf{matrix} er et rektangulært skema af tal ordnet i rækker og søjler. Et lineært ligningssystems \textbf{koefficientmatrix} indeholder kun variablenes koefficienter, mens \textbf{totalmatrixen} har en ekstra søjle med konstanterne fra ligningernes højresider. \end{definition}

\begin{example} For det lineære ligningssystem:
\[
\begin{aligned}
2 x_1+3 x_2+x_3 & =3 \\
7 x_2-4 x_3 & =10 \\
x_3 & =1
\end{aligned}
\]

er koefficientmatrixen:

\[
\left[\begin{array}{ccc}
2 & 3 & 1 \\
0 & 7 & -4 \\
0 & 0 & 1
\end{array}\right]
\]

og totalmatrixen er:

\[
\left[\begin{array}{ccc|c}
2 & 3 & 1 & 3 \\
0 & 7 & -4 & 10 \\
0 & 0 & 1 & 1
\end{array}\right]
\]

\end{example}

Den lodrette streg mellem koefficientdelen og den tilføjede søjle er valgfri og bruges til visuelt at adskille koefficienterne fra konstanterne.

\begin{definition} En matrix' \textbf{størrelse} er bestemt af antallet af rækker og søjler og skrives som \emph{rækker} $\times$ \emph{søjler}. For eksempel har koefficientmatrixen ovenfor størrelsen $3 \times 3$, og totalmatrixen har størrelsen $3 \times 4$. \end{definition}

\subsection*{Løsning af ligningssystemer med totalmatricer}

Vi viser nu, hvordan man løser et lineært ligningssystem ved hjælp af totalmatrixnotation. Fremgangsmåden formaliseres i næste afsnit, men her giver vi et eksempel. Metoden gør det lettere at løse ligningssystemer, fordi vi kun skal arbejde med matricen.

\begin{example}
    \label{ex:augmented-matrix}
Lad $r_1$ betegne den første række, $r_2$ den anden række og $r_3$ den tredje række.

\[
\left[\begin{array}{cccc}
2 & 3 & 1 & 3 \\
0 & 7 & -4 & 10 \\
0 & 0 & 1 & 1
\end{array}\right]
\]

For at forenkle systemet udfører vi \textit{elementære rækkeoperationer}. Konkret ændrer vi $r_1$ og $r_2$ således:
\[
\begin{aligned}
r_1 & \mapsto r_1 - r_3, \\
r_2 & \mapsto r_2 + 4r_3,
\end{aligned}
\]
hvilket giver:
\[
\left[\begin{array}{cccc}
2 & 3 & 0 & 2 \\
0 & 7 & 0 & 14 \\
0 & 0 & 1 & 1
\end{array}\right]
\]

Dernæst skalerer vi $r_2$ med $\frac{1}{7}$:
\[
r_2 \mapsto \frac{1}{7}r_2, \quad \left[\begin{array}{cccc}
2 & 3 & 0 & 2 \\
0 & 1 & 0 & 2 \\
0 & 0 & 1 & 1
\end{array}\right]
\]

Nu eliminerer vi $3$-tallet i den første række ved at udføre:
\[
r_1 \mapsto r_1 - 3r_2, \quad \left[\begin{array}{cccc}
2 & 0 & 0 & -4 \\
0 & 1 & 0 & 2 \\
0 & 0 & 1 & 1
\end{array}\right]
\]

Til sidst skalerer vi $r_1$ med $\frac{1}{2}$ og får:
\[
r_1 \mapsto \frac{1}{2}r_1, \quad \left[\begin{array}{cccc}
1 & 0 & 0 & -2 \\
0 & 1 & 0 & 2 \\
0 & 0 & 1 & 1
\end{array}\right]
\]

Nu kan vi igen tolke matricen som et lineært ligningssystem. Ud fra totalmatrixen får vi:
\[
\begin{aligned}
x_1 &= -2, \\
x_2 &= 2, \\
x_3 &= 1.
\end{aligned}
\]
Dette kan skrives kompakt som:
\[
(x_1, x_2, x_3) = (-2, 2, 1).
\]
Løsningen er altså den samme som den, vi fandt ved indsættelse.
\end{example}

\begin{example}
    \label{ex:inconsistent-system}
Afgør, om følgende system er konsistent:
        \[
            \begin{array}{r}
            x_2+4 x_3=2 \\
            x_1-3 x_2+2 x_3=6 \\
            x_1-2 x_2+6 x_3=9
            \end{array}
        \]
\end{example}

\begin{solution}
    Først skriver vi systemets totalmatrix:
    \[
    \left[\begin{array}{cccc}
    0 & 1 & 4 & 2 \\
    1 & -3 & 2 & 6 \\
    1 & -2 & 6 & 9
    \end{array}\right]
    \]
    
    For at forenkle matricen ombytter vi $r_1$ og $r_2$:
    \[
    r_1 \leftrightarrow r_2, \quad \left[\begin{array}{cccc}
    1 & -3 & 2 & 6 \\
    0 & 1 & 4 & 2 \\
    1 & -2 & 6 & 9
    \end{array}\right]
    \]
    
    Dernæst eliminerer vi $1$-tallet i den første søjle af $r_3$ ved at udføre:
    \[
    r_3 \mapsto r_3 - r_1, \quad \left[\begin{array}{cccc}
    1 & -3 & 2 & 6 \\
    0 & 1 & 4 & 2 \\
    0 & 1 & 4 & 3
    \end{array}\right]
    \]
    
    Derefter eliminerer vi $1$-tallet i den anden søjle af $r_3$:
    \[
    r_3 \mapsto r_3 - r_2, \quad \left[\begin{array}{cccc}
    1 & -3 & 2 & 6 \\
    0 & 1 & 4 & 2 \\
    0 & 0 & 0 & 1
    \end{array}\right]
    \]
    
    Den tredje række svarer nu til ligningen:
    \[
    0x_1 + 0x_2 + 0x_3 = 1
    \]
    som kan reduceres til:
    \[
    0 = 1
    \]
    Det er en modstrid, så systemet er \textbf{inkonsistent} og har ingen løsning. \end{solution}

\begin{example}
    Angiv en løsning til følgende system (hvis der findes en). Er den entydig?
    \[
        \begin{aligned}
        x_1+2 x_2+3 x_3 & =4 \\
        3 x_1+6 x_2+9 x_3 & =12
        \end{aligned}
    \]
\end{example}

\begin{solution}
        \[
        \left[\begin{array}{cccc}
        1 & 2 & 3 & 4 \\
        3 & 6 & 9 & 12
        \end{array}\right]
        \]
        
        Vi bemærker, at $r_2$ er et multiplum af $r_1$:
        \[
        r_2 = 3r_1
        \]
        
        For at forenkle systemet udfører vi følgende elementære rækkeoperation, som fjerner overflødig information:
        \[
        r_2 \mapsto r_2 - 3r_1,
        \]
        hvilket giver:
        \[
        \left[\begin{array}{cccc}
        1 & 2 & 3 & 4 \\
        0 & 0 & 0 & 0
        \end{array}\right]
        \]
        
        Den anden række består nu udelukkende af nuller, hvilket viser, at den ikke giver nogen ny information. Det betyder, at systemet er underbestemt og har uendeligt mange løsninger.
        
        Vi kan udtrykke $x_1$ ud fra den første ligning:
        \[
        x_1 = 4 - 2x_2 - 3x_3
        \]
        
        Vi lader $x_2$ og $x_3$ være frie variable (parametre) og sætter:
        \[
        x_2 = s, \quad x_3 = t, \quad \text{hvor } s, t \in \mathbb{R}
        \]
        
        Ved indsættelse får vi:
        \[
        \begin{aligned}
        x_1 &= 4 - 2s - 3t, \\
        x_2 &= s, \\
        x_3 &= t.
        \end{aligned}
        \]
        
        Dette kan skrives kompakt som:
        \[
        (x_1, x_2, x_3) = (4 - 2s - 3t, \ s, \ t), \quad \text{for alle } s, t \in \mathbb{R}.
        \]
        
        Systemet har uendeligt mange løsninger, der er parametriseret ved $s$ og $t$. Løsningen er derfor ikke entydig.\end{solution}

\begin{example}
    \label{ex:unique-solution}
    Vælg $h$ og $k$, så følgende system
    \[
    \begin{aligned}
    x_1-3 x_2 & =1 \\
    2 x_1+h x_2 & =k
    \end{aligned}
    \]
    har
\begin{enumerate}[label=(\roman*)]
    \item en entydig løsning,
    \item mange løsninger og
    \item ingen løsning.
\end{enumerate}

\begin{solution}
        \[
        \left[\begin{array}{ccc}
        1 & -3 & 1 \\
        2 & h & k
        \end{array}\right]
        \]
        
        For at forenkle systemet udfører vi en elementær rækkeoperation, som eliminerer $x_1$ fra den anden ligning. Konkret ændrer vi $r_2$ således:
        \[
        r_2 \mapsto r_2 - 2r_1,
        \]
        hvilket giver:
        \[
        \left[\begin{array}{ccc}
        1 & -3 & 1 \\
        0 & h + 6 & k - 2
        \end{array}\right]
        \]
        
        Nu analyserer vi det resulterende system ud fra værdien af $h + 6$.
        
        \textbf{Tilfælde 1: $h + 6 \ne 0$}
        
        Når $h + 6 \ne 0$, kan vi løse den anden ligning med hensyn til $x_2$:
        \[
        (h + 6)x_2 = k - 2 \implies x_2 = \dfrac{k - 2}{h + 6}
        \]
        Vi indsætter $x_2$ i den første ligning:
        \[
        x_1 - 3x_2 = 1 \implies x_1 = 1 + 3x_2 = 1 + 3\left(\dfrac{k - 2}{h + 6}\right)
        \]
        Vi får altså en entydig løsning for $x_1$ og $x_2$. \textbf{Systemet har en entydig løsning, når $h + 6 \ne 0$.}
        
        \newpage

        \textbf{Tilfælde 2: $h + 6 = 0$}
        
        Når $h + 6 = 0$, dvs. $h = -6$, bliver den anden ligning:
        \[
        0x_2 = k - 2
        \]
        som kan reduceres til:
        \[
        0 = k - 2
        \]
        Vi har to deltilfælde:
        
        \textbf{\textit{Deltilfælde 2a:} $k - 2 = 0$}
        
          Hvis $k = 2$, bliver ligningen $0 = 0$, som altid er sand. Den anden ligning giver ingen ny information, så systemet reduceres til:
          \[
          x_1 - 3x_2 = 1
          \]
          Her er $x_2$ en fri variabel. Vi løser med hensyn til $x_1$:
          \[
          x_1 = 1 + 3x_2
          \]
          \textbf{Systemet har uendeligt mange løsninger, når $h = -6$ og $k = 2$.}
        
        \textbf{\textit{Deltilfælde 2b:} $k - 2 \ne 0$}
        
          Hvis $k \ne 2$, bliver ligningen $0 = k - 2$, hvilket er en modstrid, da $0 \ne k - 2$. \textbf{Systemet har altså ingen løsning, når $h = -6$ og $k \ne 2$.}
                
        \textbf{Konklusion:}
        
        \begin{enumerate}[label=(\roman*)]
            \item \textbf{Ingen løsning}, når $h = -6$ og $k \ne 2$.
            \item \textbf{En entydig løsning}, når $h \ne -6$.
            \item \textbf{Uendeligt mange løsninger}, når $h = -6$ og $k = 2$.
        \end{enumerate} \end{solution}


\end{example}
\section{Elementære rækkeoperationer og echelonformer}

I \autoref{ex:augmented-matrix}--\autoref{ex:unique-solution} udførte vi en række operationer, der kaldes \textbf{elementære rækkeoperationer}. Det er grundlæggende omformninger, som forenkler ligningssystemer uden at ændre deres løsningsmængder.

De tre typer elementære rækkeoperationer er:
\begin{itemize}
    \item \textbf{Erstatning:} Erstat en række med summen af den selv og et multiplum af en anden række.
    \item \textbf{Ombytning:} Byt om på to rækker.
    \item \textbf{Skalering:} Gang alle indgange i en række med en konstant forskellig fra nul.
\end{itemize}

To matricer kaldes \textbf{rækkeækvivalente}, hvis den ene kan omdannes til den anden ved en følge af elementære rækkeoperationer.

\begin{remark}
To lineære ligningssystemer er ækvivalente (dvs. de har samme løsningsmængde), hvis deres totalmatricer er rækkeækvivalente. Elementære rækkeoperationer på en totalmatrix ændrer ikke systemets løsningsmængde. \end{remark}

Som du måske har bemærket i eksemplerne, får de matricer, der fremkommer ved disse operationer, ofte et særligt mønster -- mange nuller under bestemte indgange. Det er ikke et tilfælde, men et mål med processen. Vi kalder dette matricens \textbf{echelonform}. Ved at bringe en matrix på echelonform:

\begin{itemize}
    \item gør vi det lettere at løse lineære ligningssystemer,
    \item forenkler vi matricen systematisk, så beregningerne bliver mindre komplicerede,
    \item afslører vi vigtige egenskaber ved systemet, fx om der findes løsninger og hvor mange, eller om bestemte ligninger er afhængige.
\end{itemize}

\begin{definition}{Echelonformer}
En matrix er på \textbf{echelonform}, hvis den opfylder følgende betingelser:
\begin{enumerate}
    \item Alle nulrækker står nederst.
    \item Hver rækkes ledende indgang står i en søjle til højre for den ledende indgang i rækken ovenover.
    \item Alle indgange i en søjle under en ledende indgang er nuller.
\end{enumerate}

Ud over echelonform kan en matrix også være på \textbf{reduceret echelonform} (RREF, af engelsk \emph{reduced row echelon form}), som har følgende egenskaber:
\begin{enumerate}[resume]
    \item Den ledende indgang i hver række, der ikke er en nulrække, er 1.
    \item Hvert ledende 1-tal er den eneste indgang forskellig fra nul i sin søjle.
\end{enumerate}
\end{definition}

Vi siger, at en echelonmatrix $U$ er en echelonform af matricen $A$, hvis $U$ er rækkeækvivalent med $A$. Tilsvarende siger vi, at en reduceret echelonmatrix $U$ er den reducerede echelonform af matricen $A$, hvis $U$ er rækkeækvivalent med $A$.

Betydningen af at bringe totalmatrixen for et lineært ligningssystem på echelonform forklares af
følgende sætning.

\begin{theorem}{Eksistenssætningen}
Et lineært ligningssystem er konsistent, hvis og kun hvis en echelonform af totalmatrixen ikke har nogen række på formen
    \[
    \left[\begin{array}{llll}
    0 & \ldots & 0 & b
    \end{array}\right]
    \]
hvor $b$ er forskellig fra nul.
\end{theorem}

Vi så et eksempel på et inkonsistent system i \autoref{ex:inconsistent-system}, fordi echelonformen af totalmatrixen har en række på formen $\left[\begin{array}{lll|l} 0 & 0 & 0 & 1 \end{array}\right]$. Denne række viser, at systemet ikke har nogen løsning.

\begin{example}
    Totalmatrixen for det lineære ligningssystem, der var hovedeksemplet i det foregående afsnit,

\[
\left[\begin{array}{cccc}
2 & 3 & 1 & 3 \\
0 & 7 & -4 & 10 \\
0 & 0 & 1 & 1
\end{array}\right] , 
\] er allerede på echelonform. Da den ikke har nogen række af den form, der er nævnt i sætningen, ved vi straks,
at systemet er konsistent. De ledende indgange er $2$, $7$ og $1$, og alle indgange under dem er nuller. Matricen er dog ikke på reduceret echelonform, fordi de ledende indgange ikke alle er $1$.
\end{example}

Husk, at vi udførte en række rækkeoperationer på den foregående matrix og fik

\[
\left[\begin{array}{cccc}
1 & 0 & 0 & -2 \\
0 & 1 & 0 & 2 \\
0 & 0 & 1 & 1
\end{array}\right]
\]

som er på reduceret echelonform. Det gjorde det let at se systemets løsninger, hvilket er den største fordel ved at bringe matricen på denne form. Det fører os til følgende vigtige sætning.

\begin{theorem}{Entydighed af den reducerede echelonform}
    Enhver matrix er rækkeækvivalent med én og kun én reduceret echelonmatrix.
\end{theorem}

En pivotposition i en matrix $A$ er en placering i $A$, der svarer til et ledende 1-tal i den reducerede echelonform af $A$. En pivotsøjle er en søjle i $A$, der indeholder en pivotposition. Bemærk, at man kan få en pivot ved at skalere den ledende indgang i en række, så den bliver $1$. Pivotsøjlen er derfor den søjle, hvor den ledende indgang står.

Nu er vi klar til at præsentere algoritmen til løsning af et lineært ligningssystem ved hjælp af matricer.

\begin{custombox}{Rækkereduktionsalgoritmen}
    Her beskriver vi en algoritme, der omdanner en vilkårlig matrix til en ækvivalent
    (reduceret) echelonmatrix. Algoritmen er grundlaget for at løse lineære ligningssystemer ved hjælp af matricer.

    \begin{enumerate}
        \item Begynd med den søjle længst til venstre, der ikke er en nulsøjle. Det er en pivotsøjle, og pivotpositionen står øverst.
        \item Vælg en indgang forskellig fra nul i pivotsøjlen som pivot. Byt om nødvendigt rækker, så denne indgang flyttes til pivotpositionen.
        \item Brug erstatningsoperationer til at skabe nuller på alle positioner under pivoten.
        \item Anvend trin 1--3 på den delmatrix, der består af alle indgange under og til højre for pivotpositionen. Gentag processen, indtil der ikke er flere rækker forskellige fra nul at ændre. (Nu har vi nået en echelonform af matricen.)
        \item Begynd med pivoten længst til højre, og arbejd opad og mod venstre: Skab nuller over hver pivot ved hjælp af rækkeoperationer. Hvis en pivot ikke er 1, gøres den til 1 ved en skalering. (Dette trin giver matricens reducerede echelonform.)
    \end{enumerate}

\end{custombox}

\subsection*{Løsninger til lineære ligningssystemer}

Lad $A$ være koefficientmatrixen for et lineært ligningssystem. Pivotsøjlerne i matricen svarer til det, vi kalder \textbf{basisvariable}. Søjlerne uden pivot svarer til det, vi kalder \textbf{frie variable}.

\begin{example} Antag, at totalmatrixen for et lineært ligningssystem er reduceret til følgende form:

\[
\left[\begin{array}{lllll}
\blacksquare & * & * & * & * \\
0 & \blacksquare & * & * & * \\
0 & 0 & 0 & \blacksquare & * \\
0 & 0 & 0 & 0 & 0
\end{array}\right]
\]

hvor $\blacksquare$ betegner et vilkårligt tal forskelligt fra nul, og $*$ betegner et vilkårligt tal (inklusive 0 ). Systemets basisvariable er $x_1, x_2$ og $x_4$. Den eneste frie variabel er $x_3$.

\end{example}


\begin{theorem}{Entydighedssætningen}
    Hvis et lineært ligningssystem er konsistent, indeholder løsningsmængden enten
    \begin{enumerate}[label=(\roman*)]
        \item en entydig løsning, når der ikke er nogen frie variable, eller
        \item uendeligt mange løsninger, når der er mindst én fri variabel.
    \end{enumerate}
\end{theorem}

\begin{example} Antag, at følgende matrix er totalmatrixen for et lineært ligningssystem i variablene $x_1, x_2$ og $x_3$. Rækkereducér matricen til echelonform for at afgøre, om systemet er konsistent.

\[
\left[\begin{array}{cccc}
1 & 2 & 3 & 4 \\
5 & 6 & 7 & 8 \\
9 & 10 & 11 & 12
\end{array}\right]
\]
Hvis systemet er konsistent, så find den reducerede echelonform, og skriv løsningsmængden med de frie variable som parametre.
\end{example}

\begin{solution} For at forenkle systemet udfører vi \textit{elementære rækkeoperationer}.
        
    \textbf{Trin 1: Eliminér \( 5 \)-tallet i \( r_2 \)}
        
        Vi ændrer \( r_2 \) således:
        \[
        r_2 \mapsto r_2 - 5r_1,
        \]
        hvilket giver:
        \[
        \left[\begin{array}{cccc}
        1 & 2 & 3 & 4 \\
        0 & -4 & -8 & -12 \\
        9 & 10 & 11 & 12
        \end{array}\right]
        \]
        
    \textbf{Trin 2: Eliminér \( 9 \)-tallet i \( r_3 \)}
        
        Vi ændrer \( r_3 \) således:
        \[
        r_3 \mapsto r_3 - 9r_1,
        \]
        hvilket giver:
        \[
        \left[\begin{array}{ccc|c}
        1 & 2 & 3 & 4 \\
        0 & -4 & -8 & -12 \\
        0 & -8 & -16 & -24
        \end{array}\right]
        \]
        
    \textbf{Trin 3: Eliminér \( -8 \)-tallet i \( r_3 \)}
        
        Vi ændrer \( r_3 \) igen for at eliminere indgangen i den anden søjle:
        \[
        r_3 \mapsto r_3 - 2r_2,
        \]
        hvilket giver:
        \[
        \left[\begin{array}{cccc}
        1 & 2 & 3 & 4 \\
        0 & -4 & -8 & -12 \\
        0 & 0 & 0 & 0
        \end{array}\right]
        \]
        
    \newpage

    \textbf{Trin 4: Skalér \( r_2 \) for at få et ledende 1-tal}
        
        Vi skalerer \( r_2 \) med \( -\dfrac{1}{4} \):
        \[
        r_2 \mapsto -\dfrac{1}{4} r_2,
        \]
        hvilket giver:
        \[
        \left[\begin{array}{ccc|c}
        1 & 2 & 3 & 4 \\
        0 & 1 & 2 & 3 \\
        0 & 0 & 0 & 0
        \end{array}\right]
        \]
        
    \textbf{Trin 5: Eliminér \( 2 \)-tallet i \( r_1 \)}
        
        Vi ændrer \( r_1 \) således:
        \[
        r_1 \mapsto r_1 - 2r_2,
        \]
        hvilket giver:
        \[
        \left[\begin{array}{ccc|c}
        1 & 0 & -1 & -2 \\
        0 & 1 & 2 & 3 \\
        0 & 0 & 0 & 0
        \end{array}\right]
        \]
        
    Nu er matricen på \textit{reduceret echelonform}.
        
    \textbf{Trin 6: Tolk matricen som et lineært ligningssystem}
        
        Ud fra totalmatrixen har vi:
        \[
        \begin{aligned}
        r_1 &: \quad x_1 - x_3 = -2, \\
        r_2 &: \quad x_2 + 2x_3 = 3, \\
        r_3 &: \quad 0 = 0.
        \end{aligned}
        \]
        
    Da den tredje række svarer til ligningen \( 0 = 0 \), som altid er sand, er systemet \textbf{konsistent}.
        
    \textbf{Trin 7: Udtryk løsningen ved hjælp af en fri variabel}
        \[
        \begin{aligned}
        x_1 &= -2 + x_3, \\
        x_2 &= 3 - 2x_3, \\
        x_3 &= x_3.
        \end{aligned}
        \]      
\end{solution}
